% Generated by mceg equations; edit the contract, then regenerate.
% Requires amsmath. Authored math is retained; TeX is not executed here.
% Definitions and hashes do not establish scientific validity.

% equation-sha256: 6c7b52cbfdac3da24c7a52e989f9a6db100292facd652dd3069fabb14dc40a0b
\paragraph{Equation statistical\_shells (revision 1)}
The added shells distribute population statistically over orbital states and sum each physical transition into its destination group.
\begin{equation}
\begin{aligned}
g_n=2n^2,\quad p_{n\ell}=\frac{2\ell+1}{n^2},\\\overline{A}_{nj}=\sum_{\ell=0}^{n-1}p_{n\ell}\sum_{k\in j}A_{n\ell,k},\qquad\overline{\sigma}_n(\nu)=\sum_{\ell=0}^{n-1}p_{n\ell}\sigma_{n\ell}(\nu)
\end{aligned}
\label{eq:escape.atom_5f_convergence:statistical_5f_shells:1}
\end{equation}
\noindent\textbf{Scope:} Resolved 1s,2s,2p,3s,3p,3d; collapsed shells 4 through N, with 3<=N<=20.

\begin{description}
\item[$n$] Principal quantum number (1; ID: principal).
\item[$\ell$] Orbital quantum number (1; ID: orbital).
\item[$g_n,p_{n\ell}$] State degeneracy and orbital population fraction (1; ID: weight).
\item[$A,\overline A$] Spontaneous electric-dipole rate (s\textasciicircum{}-1; ID: einstein\_a).
\item[$\sigma,\overline\sigma$] Photoionization cross section (cm\textasciicircum{}2; ID: cross\_section).
\item[$\nu$] Ordinary photon frequency (Hz; ID: frequency).
\item[$N$] Largest included principal shell (1; ID: maximum\_shell).
\end{description}

\noindent\textbf{Assumptions:}
\begin{itemize}
\item Statistical orbital populations in the collapsed shells are an approximation, not a new resolved high-l kinetic model.
\item Low-level energies and transition rates retain the selected original data.
\item Every retained shell has explicit recombination, ionization and cascade paths; omitted shells have no artificial top-off source.
\end{itemize}

% equation-sha256: b09bb84c99419c1f81e335017fff23cf849deba6cf8ad11ce2a04811249e6373
\paragraph{Equation radial\_dipoles (revision 1)}
Hydrogenic electric-dipole radial integrals; exact rational polynomial integration prevents cancellation in the static construction.
\begin{equation}
\begin{aligned}
\widetilde R_{n\ell}(u)=\frac{2}{n^2}\sqrt{\frac{(n-\ell-1)!}{(n+\ell)!}}e^{-u/n}\left(\frac{2u}{n}\right)^\ell L_{n-\ell-1}^{2\ell+1}\!\left(\frac{2u}{n}\right),\\ d_{n\ell,n^\prime\ell^\prime}=a_{\rm H}\int_0^\infty u^3\widetilde R_{n\ell}(u)\widetilde R_{n^\prime\ell^\prime}(u)\,du,\\ A_{n\ell,n^\prime\ell^\prime}=\frac{64\pi^4 e^2\nu^3}{3hc^3}\frac{\max(\ell,\ell^\prime)}{2\ell+1}d^2_{n\ell,n^\prime\ell^\prime}
\end{aligned}
\label{eq:escape.atom_5f_convergence:radial_5f_dipoles:1}
\end{equation}
\noindent\textbf{Scope:} 1<=n,nprime<=20; n>nprime and |l-lprime|=1 for a downward radiative transition.

\begin{description}
\item[$u=r/a_{\rm H}$] Dimensionless radius (1; ID: radius).
\item[$d$] Radial dipole integral (cm; ID: radial).
\item[$A$] Spontaneous rate averaged over upper magnetic states (s\textasciicircum{}-1; ID: einstein\_a).
\item[$\nu$] Finite-state transition frequency (Hz; ID: frequency).
\item[$n$] Upper principal quantum number (1; ID: n\_upper).
\item[$\ell$] Upper orbital quantum number (1; ID: l\_upper).
\item[$n^\prime$] Lower principal quantum number (1; ID: n\_lower).
\item[$\ell^\prime$] Lower orbital quantum number (1; ID: l\_lower).
\item[$a_{\rm H}$] Protium reduced-mass Bohr radius used by the bound-bound construction (cm; ID: bohr\_radius).
\end{description}

\noindent\textbf{Assumptions:}
\begin{itemize}
\item The Bohr radius uses the declared protium reduced-mass correction.
\item Polynomial summation is exact rational arithmetic; the final normalization and returned rates use float64.
\item This static atomic construction is independent of differentiable gas parameters.
\end{itemize}

% equation-sha256: bbd1bf2e3471abef0371c60714091712293206dbf51bcfc97ccd160674c8de88
\paragraph{Equation coulomb\_bound\_free (revision 1)}
Coulomb-wave cross sections recombine the normalized radial recurrence in logarithmic form.
\begin{equation}
\begin{aligned}
E=\frac{h\nu}{\chi_{n\ell}},\quad t^2=\max(E-1,0),\quad \mathcal E=1+t^2,\quad k=t/n,\\
\log G_* =\tfrac12\log(\pi/2)+\log(8n)+n\log(4n)-\log\Gamma(2n)\\
\hspace{7mm}-2n\frac{\arctan t}{t}-\tfrac12\log(1-\exp[-2\pi/k])-(n+2)\log\mathcal E,\\
\log D_{n\ell,\ell'}^2=\log\frac{(n+\ell)!}{(n-\ell-1)!}
 +\sum_{s=1}^{\ell'}\log(1+s^2k^2)+2(\ell-n)\log(2n)\\
\hspace{7mm}+2[\log G_*+p_\ell^\pm\log\mathcal E+\log|H_\ell^\pm|],\\
p_\ell^+=n-1-\ell,\quad p_\ell^-=n-\ell,\\
\sigma_{n\ell}(\nu)=\frac{4\pi\alpha a_0^2n^2}{3}
 \sum_{\ell'=\ell\pm1\ge0}\frac{\max(\ell,\ell')}{2\ell+1}\,\mathcal E\,\exp(\log D_{n\ell,\ell'}^2)
\end{aligned}
\label{eq:escape.atom_5f_convergence:coulomb_5f_bound_5f_free:1}
\end{equation}
\noindent\textbf{Scope:} 1<=n<=20; cross section zero for E<1. Plus/minus denotes continuum lprime=l+1/l-1. Each finite state uses its actual threshold.

\begin{description}
\item[$n$] Principal quantum number (1; ID: principal).
\item[$\ell$] Bound-state orbital quantum number (1; ID: orbital).
\item[$E$] Photon energy divided by the actual threshold (1; ID: energy).
\item[$G_*,H^\pm,D$] Coulomb radial normalization, normalized recurrence and dipoles (1; ID: radial).
\item[$\sigma_{n\ell}$] Orbital photoionization cross section (cm\textasciicircum{}2; ID: cross\_section).
\item[$a_0$] Bohr radius 5.29177210903e-9 cm used in the bound-free normalization (cm; ID: bohr\_radius).
\end{description}

\noindent\textbf{Assumptions:}
\begin{itemize}
\item At threshold arctan(t)/t tends to one; a small-t series preserves the derivative.
\item Computing the logarithm of the product avoids squaring tiny dipoles before compensating factors are applied.
\item Both photoionization and the inverse Milne process use this same cross section.
\item The support limit n<=20 is explicit; successful tests at sampled energies do not prove an unlimited high-n domain.
\end{itemize}

% equation-sha256: 13b01dd08fb23dad58e64a975e989a47d811317651cc18de0cf86be28e16c893
\paragraph{Equation normalized\_coulomb\_recurrence (revision 1)}
Energy-normalized Brocklehurst recurrence, algebraically equivalent to the earlier unscaled radial recurrence.
\begin{equation}
\begin{aligned}
G_q^+=G_*\mathcal E^{n-1-q}H_q^+,\quad G_q^-=G_*\mathcal E^{n-q}H_q^-,\\
H^+_{n-1}=1,\quad H^+_{n-2}=n(2n-1),\\
H^-_{n-1}=\frac1{2n},\quad H^-_{n-2}=(2n-1)\left(\frac4{\mathcal E}+n-1\right)H^-_{n-1},\\
A_q^+=\frac{4n^2-4(q+2)^2}{\mathcal E}+(q+2)(2q+3),\\
B_q^+=\frac{4n^2[n^2-(q+2)^2][1+(q+3)^2k^2]}{\mathcal E^2},\\
A_q^-=\frac{4n^2-4(q+1)^2}{\mathcal E}+(q+1)(2q+3),\\
B_q^-=\frac{4n^2[n^2-(q+2)^2][1+(q+1)^2k^2]}{\mathcal E^2},\\
H_q^\pm=A_q^\pm H_{q+1}^\pm-B_q^\pm H_{q+2}^\pm
\end{aligned}
\label{eq:escape.atom_5f_convergence:normalized_5f_coulomb_5f_recurrence:1}
\end{equation}
\noindent\textbf{Scope:} Plus runs q=n-3,...,0; minus runs q=n-3,...,1. Initial values are used only when their indices exist. The noncontributing n=1 minus seed is never evaluated as a dipole.

\begin{description}
\item[$n$] Principal quantum number (1; ID: principal).
\item[$E$] Photon energy divided by threshold; the clipped energy is max(E,1) (1; ID: energy).
\item[$k$] Dimensionless continuum momentum (1; ID: continuum).
\item[$H_q^\pm$] Energy-normalized radial factors (1; ID: radial).
\end{description}

\noindent\textbf{Assumptions:}
\begin{itemize}
\item Each downward step removes one power of the clipped photon energy before forming the recurrence.
\item The algorithm is a provider-derived rescaling of the inspected Cloudy/Brocklehurst recurrence, not an empirical population adjustment.
\end{itemize}

% equation-sha256: 022388d9dab1fa55ffdb5170110c0b8d9f123eb7c621395ac556c8f2b87132aa
\paragraph{Equation higher\_shell\_ionization (revision 1)}
VS80 Eq. 8 supplies the added-shell electron ionization rate; its three-body inverse uses the same state energy and statistical weight.
\begin{equation}
\begin{aligned}
u=\frac{k_{\rm B}T}{1\,\mathrm{eV}},\quad\epsilon_n=\frac{\chi_n}{k_{\rm B}T},\\q_n^{\rm ion}=\frac{9.56\times10^{-6}e^{-\epsilon_n}}{u^{3/2}[\epsilon_n^{2.33}+4.38\epsilon_n^{1.72}+1.32\epsilon_n]},\\q_n^{\rm 3b}=q_n^{\rm ion}\frac{g_n}{2}\left(\frac{h^2}{2\pi m_e k_{\rm B}T}\right)^{3/2}e^{\chi_n/(k_{\rm B}T)}
\end{aligned}
\label{eq:escape.atom_5f_convergence:higher_5f_shell_5f_ionization:1}
\end{equation}
\noindent\textbf{Scope:} Only collapsed n=6,...,20 use this extension. The retained n<=5 collision/ionization prescriptions remain distinct.

\begin{description}
\item[$T$] Gas temperature (K; ID: temperature).
\item[$\chi_n$] Actual finite-state ionization threshold (erg; ID: binding).
\item[$q_n^{\rm ion}$] Electron-impact ionization coefficient (cm\textasciicircum{}3 s\textasciicircum{}-1; ID: rate).
\item[$q_n^{\rm 3b}$] Three-body recombination coefficient (cm\textasciicircum{}6 s\textasciicircum{}-1; ID: inverse).
\end{description}

\noindent\textbf{Assumptions:}
\begin{itemize}
\item Maxwellian electrons; the VS80 fit is an adopted approximation.
\item The retained low-level data restrict the shared gas-temperature domain.
\end{itemize}

% equation-sha256: f10a6b3985bb3bf8ef6c50f33864e5d9629a0ac6a6027b056489d7e746796620
\paragraph{Equation finite\_charge\_root (revision 1)}
The retained charge-conserving statistical-equilibrium root is solved for the enlarged rate matrix and differentiated implicitly.
\begin{equation}
\begin{aligned}
P_{ij}=R_{ij}+n_e q_{ij},\quad
 M=\operatorname{diag}\!\left(\sum_jP_{ij}+\Gamma_i+n_e q_i^{\rm ion}\right)-P^{\mathsf T},\\
s_i=n_e\alpha_i+n_e^2q_i^{\rm 3b},\quad b=M^{-1}s,\\
n_e=n_p=n_{\rm H}e^z,\quad n_i=n_e b_i,\quad F(z;\theta)=z+\log(1+\mathbf1^{\mathsf T}b)=0,\\
\frac{dz}{d\theta}=-\frac{\partial_\theta F}{\partial_z F}
\end{aligned}
\label{eq:escape.atom_5f_convergence:finite_5f_charge_5f_root:1}
\end{equation}
\noindent\textbf{Scope:} Pure H, nonnegative finite rates with a positive photoionization rate, nonsingular bound-state matrix, isolated scalar charge root. Gas and radiation inputs are prescribed in this local comparison.

\begin{description}
\item[$n_{\rm H}$] Hydrogen nucleus number density (cm\textasciicircum{}-3; ID: density).
\item[$T$] Prescribed gas temperature (K; ID: temperature).
\item[$n_e=n_p$] Electron and proton density (cm\textasciicircum{}-3; ID: electron).
\item[$n_i$] Bound-state populations (cm\textasciicircum{}-3; ID: population).
\item[$b_i$] Bound-to-proton population ratio (1; ID: ratio).
\item[$P,R,\Gamma,M,s$] Transition rates, loss matrix and source per proton (s\textasciicircum{}-1; ID: rate).
\item[$q_{ij},q_i^{\rm ion},\alpha_i$] Two-body rate coefficients (cm\textasciicircum{}3 s\textasciicircum{}-1; ID: coefficient).
\item[$q_i^{\rm 3b}$] Three-body coefficient (cm\textasciicircum{}6 s\textasciicircum{}-1; ID: three\_body).
\item[$z,F$] Log ionized fraction and conservation residual (1; ID: root).
\end{description}

\noindent\textbf{Assumptions:}
\begin{itemize}
\item Orbital mixing coefficients remain inside the density-dependent rate callback.
\item The implicit derivative follows the converged equations, including rate feedback, rather than differentiating the bisection iterations.
\item The gas-temperature limits of the retained low-level atomic data remain applicable.
\end{itemize}

% equation-sha256: f5c9364243ceecec01d7958d52e8a4f2a56c0157e55df252ff2915c7c2691925
\paragraph{Equation shell\_change (revision 1)}
A sampled relative shell-size criterion with the smaller atom in the denominator.
\begin{equation}
\begin{aligned}
\Delta_s(N,N_{\rm ref})=\left|\frac{n_s(N_{\rm ref})}{n_s(N)}-1\right|,\qquad\max_{s\in\{2,3\}}\Delta_s(15,20)\le0.05
\end{aligned}
\label{eq:escape.atom_5f_convergence:shell_5f_change:1}
\end{equation}
\noindent\textbf{Scope:} The seven registered temperature/density/radiation cases; N=15 and Nref=20. The separate baseline measurement retains 48 points per threshold interval.

\begin{description}
\item[$n_s$] Population summed over resolved sublevels of shell s (cm\textasciicircum{}-3; ID: population).
\item[$\Delta_s$] Absolute relative change on increasing the largest included shell (1; ID: change).
\end{description}

\noindent\textbf{Assumptions:}
\begin{itemize}
\item Each population solve must meet its numerical convergence criteria.
\item Numerical quadrature and derivative tests are separate prerequisites.
\item This finite comparison cannot establish infinite-atom convergence or physical atmospheric accuracy.
\end{itemize}
